% GENERATED FILE. Edit canonical lessons, then run npm run generate:books.
% canonical-source-hash: fnv1a64:6bbc0fad754171ef
% canonical-lessons: PT-C438-subida, PT-C438-descida, PT-C438-ladeira, PT-C438-ladear, PT-C438-desembocar

\chapter{Subir e descer}
\label{ch:pt-subir-e-descer}
\hlchaptermodality{438}

\begin{chapteropening}
\textbf{By the end of this chapter:} \emph{I can describe a route: the climb, the way down, a steep street, running alongside and coming out somewhere.}

Complete the last lesson of chapter 438: I can describe a route: the climb, the way down, a steep street, running alongside and coming out somewhere.
\end{chapteropening}

\section[a subida]{a subida — the climb, the way up}

\label{lesson:PT-C438-subida}



\begin{quote}
Before the new one: say the Portuguese for the narrow lane, then the Portuguese for the fork.
\end{quote}

\subsection*{You'll want to know: a subida}
\textbf{a subida} — \textquotedblleft{}the climb, the way up\textquotedblright{}.

\begin{quote}
A subida até ao castelo é cansativa.
\end{quote}

\begin{quote}
\emph{The climb up to the castle is tiring.}
\end{quote}

Also a rise: \emph{uma subida de preços}.

\begin{cousinweb}
\emph{subida} is built on \emph{subir} 'to go up' (EtymDB), which is inherited from Latin \emph{subīre} 'to go under, to come up'.
\end{cousinweb}

\begin{grammarlens}[title={What you've built}]
One more word for telling, explaining and describing at B1.
\end{grammarlens}

\subsection*{Your turn}
\begin{itemize}
\raggedright
  \item \emph{Say it:} \emph{a subida}
  \item \emph{Say it:} \emph{a subida}, once more
  \item \emph{Say it:} a sentence of your own with \emph{a subida}
  \item \emph{Recall:} say the Portuguese for the narrow lane, then the Portuguese for the fork, then say \emph{a subida} again
\end{itemize}

\begin{tcolorbox}[breakable,colback=teal!4,colframe=teal!35!black,title={Before you move on}]
What does a subida mean? (\textquotedblleft{}The climb, the way up\textquotedblright{}.) Say it once more.
\end{tcolorbox}
\section[a descida]{a descida — the way down, the descent}

\label{lesson:PT-C438-descida}



\begin{quote}
Before the new one: say the Portuguese for the fork, then the Portuguese for the climb.
\end{quote}

\subsection*{You'll want to know: a descida}
\textbf{a descida} — \textquotedblleft{}the way down, the descent\textquotedblright{}.

\begin{quote}
Na descida, vamos devagar por causa das pedras.
\end{quote}

\begin{quote}
\emph{On the way down, we go slowly because of the stones.}
\end{quote}

Also a fall: \emph{uma descida da temperatura}.

\begin{cousinweb}
\emph{descida} is built on \emph{descer} 'to go down', which is inherited from Latin \emph{dēscendere}, on \emph{scandere} 'to climb' (Lewis \& Short).
\end{cousinweb}

\begin{grammarlens}[title={What you've built}]
One more word for telling, explaining and describing at B1.
\end{grammarlens}

\subsection*{Your turn}
\begin{itemize}
\raggedright
  \item \emph{Say it:} \emph{a descida}
  \item \emph{Say it:} \emph{a descida}, once more
  \item \emph{Say it:} a sentence of your own with \emph{a descida}
  \item \emph{Recall:} say the Portuguese for the fork, then the Portuguese for the climb, then say \emph{a descida} again
\end{itemize}

\begin{tcolorbox}[breakable,colback=teal!4,colframe=teal!35!black,title={Before you move on}]
What does a descida mean? (\textquotedblleft{}The way down, the descent\textquotedblright{}.) Say it once more.
\end{tcolorbox}
\section[a ladeira]{a ladeira — the steep street, the slope}

\label{lesson:PT-C438-ladeira}



\begin{quote}
Before the new one: say the Portuguese for the climb, then the Portuguese for the way down.
\end{quote}

\subsection*{You'll want to know: a ladeira}
\textbf{a ladeira} — \textquotedblleft{}the steep street, the slope\textquotedblright{}.

\begin{quote}
Moramos no cimo de uma ladeira.
\end{quote}

\begin{quote}
\emph{We live at the top of a steep street.}
\end{quote}

Old hill towns are full of \emph{ladeiras}.

\begin{cousinweb}
\emph{ladeira} is built on \emph{lado} 'side' (EtymDB), which is inherited from Latin \emph{latus} 'side, flank': the side of a hill.
\end{cousinweb}

\begin{grammarlens}[title={What you've built}]
One more word for telling, explaining and describing at B1.
\end{grammarlens}

\subsection*{Your turn}
\begin{itemize}
\raggedright
  \item \emph{Say it:} \emph{a ladeira}
  \item \emph{Say it:} \emph{a ladeira}, once more
  \item \emph{Say it:} a sentence of your own with \emph{a ladeira}
  \item \emph{Recall:} say the Portuguese for the climb, then the Portuguese for the way down, then say \emph{a ladeira} again
\end{itemize}

\begin{tcolorbox}[breakable,colback=teal!4,colframe=teal!35!black,title={Before you move on}]
What does a ladeira mean? (\textquotedblleft{}The steep street, the slope\textquotedblright{}.) Say it once more.
\end{tcolorbox}
\section[ladear]{ladear — to run alongside, to flank}

\label{lesson:PT-C438-ladear}



\begin{quote}
Before the new one: say the Portuguese for the way down, then the Portuguese for the steep street.
\end{quote}

\subsection*{You'll want to know: ladear}
\textbf{ladear} — \textquotedblleft{}to run alongside, to flank\textquotedblright{}.

\begin{quote}
Árvores altas ladeiam a avenida.
\end{quote}

\begin{quote}
\emph{Tall trees line the avenue.}
\end{quote}

Also to avoid something: \emph{ladear uma pergunta}.

\begin{cousinweb}
\emph{ladear} is built on \emph{lado} 'side', which is inherited from Latin \emph{latus} 'side, flank' (EtymDB).
\end{cousinweb}

\begin{grammarlens}[title={What you've built}]
One more word for telling, explaining and describing at B1.
\end{grammarlens}

\subsection*{Your turn}
\begin{itemize}
\raggedright
  \item \emph{Say it:} \emph{ladear}
  \item \emph{Say it:} \emph{ladear}, once more
  \item \emph{Say it:} a sentence of your own with \emph{ladear}
  \item \emph{Recall:} say the Portuguese for the way down, then the Portuguese for the steep street, then say \emph{ladear} again
\end{itemize}

\begin{tcolorbox}[breakable,colback=teal!4,colframe=teal!35!black,title={Before you move on}]
What does ladear mean? (\textquotedblleft{}To run alongside, to flank\textquotedblright{}.) Say it once more.
\end{tcolorbox}
\section[desembocar]{desembocar — to come out (into); to flow (into)}

\label{lesson:PT-C438-desembocar}



\begin{quote}
Before the new one: say the Portuguese for the steep street, then the Portuguese for to run alongside.
\end{quote}

\subsection*{You'll want to know: desembocar}
\textbf{desembocar} — \textquotedblleft{}to come out (into); to flow (into)\textquotedblright{}.

\begin{quote}
Esta rua desemboca na praça da igreja.
\end{quote}

\begin{quote}
\emph{This street comes out into the church square.}
\end{quote}

A river \emph{desemboca} in the sea.

\begin{cousinweb}
\emph{desembocar} is built from \emph{des-}, \emph{em} and \emph{boca} 'mouth', inherited from Latin \emph{bucca} 'cheek, mouth' (Lewis \& Short): to come out of a mouth.
\end{cousinweb}

\begin{grammarlens}[title={What you've built}]
One more word for telling, explaining and describing at B1.
\end{grammarlens}

\subsection*{Your turn}
\begin{itemize}
\raggedright
  \item \emph{Say it:} \emph{desembocar}
  \item \emph{Say it:} \emph{desembocar}, once more
  \item \emph{Say it:} a sentence of your own with \emph{desembocar}
  \item \emph{Recall:} say the Portuguese for the steep street, then the Portuguese for to run alongside, then say \emph{desembocar} again
\end{itemize}

\begin{tcolorbox}[breakable,colback=teal!4,colframe=teal!35!black,title={Before you move on}]
What does desembocar mean? (\textquotedblleft{}To come out (into); to flow (into)\textquotedblright{}.) Say it once more.
\end{tcolorbox}
